\scanpage{124}\mspage{69}
On trouve \deleted{\unsure{pour} \unsure{$i$}}
\diagram{\begin{tikzcd}[ampersand replacement=\&,row sep=small,column sep=small]
 {} \&x\arrow[d,no head]\&{}\& {}\&\text{chemin vide}\\
 i_{C(X)}(x,c)=(x,c,1_t):\&f_X(x)\arrow[r,"c"]\&t\arrow[r,"{1_t}"{name=one}]\&t\arrow[from=one,to=1-5]\&{}
\end{tikzcd}}
\diagram{\begin{tikzcd}[ampersand replacement=\&,row sep=small,column sep=small]
 {}\&x\arrow[d]\&{}\&{}\\
 C(i_X)(x,c)=(x,1_{f_X(x)},c):\&f_X(x)\arrow[r,no head,"{1_{\unsure{f(x)|f_X(x)}}}"{name=empty}]\&f_X(x)\arrow[r,"c"]\&t\\
 {}\&{}\&\text{chemin vide}\arrow[from=empty,to=3-3]\&{}
\end{tikzcd}}

Mais on a le choix entre deux théories des chemins, pour construire $C(X)$\,: soit avec les $\text{\uline{ch}}_{\infty}(S)$, soit avec les $\text{\uline{Ch}}(S)$, ``Ça donne les mêmes résultats'', mais suivant les situations l'un est plus commode que l'autre. Je vais travailler avec \unsure{``les''|les} $\text{\uline{Ch}}(S)$, auquel cas les $1_s$ sont \deleted{des} unités à gauche et à droite, \deleted{\unsure{on}} et on a une loi de composition pour les chemins. Cette loi nous permet de définir
\[
 \begin{aligned}
 C^2(X)&\xrightarrow{\unsure{r|\gamma}}C(X)\\
 (x,c,c')&\longmapsto(x,c'\circ c),
 \end{aligned}
\]
et les composés de $i_{C(X)}$ et de $\unsure{C_{i(X)}|C(i_X)}$ avec cette flèche sont $\mathrm{id}_{C(X)}$ dans les deux cas
\diagram{\begin{tikzcd}[ampersand replacement=\&]
 C(X)\arrow[r,shift left]\arrow[r,shift right]\arrow[rr,bend right=27,"\mathrm{id}"']\&
 C^2(X)\arrow[r,"r"]\& C(X)\,.
\end{tikzcd}}
Donc pour montrer que $i_{C(X)}=C(i_X)$ dans $\mathrm{Cat}_{/S}W_S^{d\, -1}$, il suffit de prouver que $r$ est dans $W_S^d$. Mais on sait que le foncteur $C$ \uline{transforme} $W_S$ en $W_S^d$, donc $C(i_X)$ est dans $W_S^d$, et \inserted{$r\circ C(i_X)=\mathrm{id}$,} \deleted{$\unsure{C(i_X)}$} $r$ est dans $W_S^d$, OK !
% archive sheet number: 123
